\chapter{Yogasana: The Syllabus Postures}
\label{ch:asanas}

\emph{Which postures must the candidate perform, and what must the candidate be able to say about each?}
The syllabus lists standing, sitting, prone and supine asanas. For each one the candidate should know the name, the starting position, the breath, the hold, the way out, one benefit and one caution. This chapter gives them in the syllabus groups. The theory of asana, its meaning and principles, is in the introduction to yogasana (see \cref{ch:asana-intro}).

\begin{unverified}
This chapter is written from a notes block whose fact check is still pending. The postures of the Common Yoga Protocol follow that booklet. The others come from the sources named beside them: Swami Satyananda's \emph{Asana Pranayama Mudra Bandha} (a book the Board names), or a government secondary source (NIOS, CBSE), or the workshop slides; each is marked. A certified teacher must review every posture, in text and in figure, before print.
\end{unverified}

\section{How to approach every asana}

The same frame applies to every posture. Move into it and out of it slowly, with the breath. Hold it for the time given, and breathe normally during the hold unless told otherwise. Stay within your capacity, and rest afterwards in a relaxation posture. The hold times in this chapter are those of the common protocol where it gives them.

The benefits given below are effects the tradition credits to each posture. They are reasons to practise, not medical claims, and no posture is a treatment. The cautions are the other half of every posture and are taught with it. A caution is a reason to modify or skip, never a diagnosis: a student with an injury, a pregnancy, high blood pressure or heart disease gets a gentle form or another posture, and a doctor's advice where in doubt. A student who cannot reach the floor uses the modified form given in the text (the knee instead of the foot, a cushion under the knees, the hands at the waist) before anything else is tried.

The common protocol covers only part of the syllabus list. Learn its postures first, because the exam and public classes are built on them.

The syllabus groups the postures by the starting position:
\begin{itemize}
\item \textbf{Standing:} tadasana, vrikshasana, ardha chakrasana, padahastasana, kati chakrasana, trikonasana.
\item \textbf{Sitting:} dandasana, sukhasana, padmasana, vajrasana, bhadrasana, mandukasana, ushtrasana, shashankasana, uttana mandukasana, paschimottanasana, purvottanasana, vakrasana, gomukhasana.
\item \textbf{Prone (on the stomach):} bhujangasana, shalabhasana, makarasana.
\item \textbf{Supine (on the back):} pavanamuktasana, uttanapadasana, ardha halasana, setubandhasana, vipareetakarani, saral matsyasana, shavasana.
\end{itemize}

\section{Standing postures}

\subsection*{Tadasana (palm tree)}
Stand with the feet about two inches apart. Interlock the fingers and turn the palms outwards. Breathe in and raise the arms overhead in line with the shoulders, then raise the heels and balance on the toes for about 10 to 15 seconds. Breathe out, lower the heels and release the hands. Spread the weight evenly over both feet, keep the thighs firm and the chest lifted. It is credited with stability and better posture.
\textit{Caution:} avoid the heel raise with acute heart problems, varicose veins or vertigo.
\ycbfig{ycb1-fig-tadasan}

\subsection*{Vrikshasana (tree)}
Stand with the feet close and fix the eyes on a point ahead. Breathing out, bend the right leg and place the foot on the inside of the left thigh, heel towards the perineum. Breathing in, raise the arms and join the palms. Hold for 10 to 30 seconds with normal breathing, come down, and repeat on the other side. The hands may stay joined at the chest instead of going overhead, and the standing knee may stay soft. It is credited with balance and toned legs.
\textit{Caution:} avoid with arthritis, vertigo or obesity; with a knee or ankle injury do not take the foot to the thigh.
\ycbfig{ycb1-fig-vrikshasan}

\subsection*{Ardha chakrasana (half wheel)}
Support the waist with the fingers together, pointing forward or down. Drop the head back, and breathing in, bend back from the lower back. Hold for 10 to 30 seconds with normal breathing, then come up on an in-breath. The workshop slides also call it hastottanasana, but the guidebook uses that name for an arms-up posture, so the alias is not confirmed.
\textit{Caution:} avoid with vertigo or giddiness; with high blood pressure bend with care.
\ycbfig{ycb1-fig-ardhachakrasan}

\subsection*{Padahastasana (hands to feet)}
Breathing in, raise the arms and stretch up from the waist. Breathing out, bend forward until the trunk is level, then bring the palms to the floor as far as the back allows. Hold for 10 to 30 seconds. Breathing in, come up with the arms overhead; breathing out, lower the arms. A stiff back bends only as far as it can, and the knees may bend slightly when the backs of the legs are tight; the aim is the stretch, not touching the floor.
\textit{Caution:} avoid with heart or back problems, spine or disc disorders, abdominal inflammation, hernia, ulcers, high myopia, vertigo and pregnancy.
\ycbfig{ycb1-fig-padahastasan}

\subsection*{Kati chakrasana (waist wheel), the standing twist}
\textit{Source: Satyananda, a Board-named book.} Stand with the feet about half a metre apart and breathe in, raising the arms to shoulder level. Breathe out and twist to the left: the right hand goes to the left shoulder and the left arm wraps round the back to the right side of the waist. Look over the left shoulder with the neck straight and the feet fixed. Hold the breath out for about two seconds and deepen the twist, then breathe in and return. Repeat on the other side; that is one round. Do 5 to 10 rounds, relaxed and without jerking. It is credited with toning the waist, back and hips. A side bend called ardha kati chakrasana is taught elsewhere; it is a different posture.
\textit{Caution:} the source gives none; take care with back or spine problems (our advice).
\ycbfig{ycb1-fig-katichakrasan}

\subsection*{Trikonasana (triangle)}
Stand with the feet comfortably apart, the front foot turned out and the back foot slightly in, and the arms out to the sides at shoulder height. Breathing out, bend to the right and place the right hand just behind the right foot; the left arm goes straight up in line with the right, palm forward. Turn the head and look at the tip of the left middle finger. Hold for 10 to 30 seconds, come up on an in-breath, and repeat on the left. If the foot is out of reach, rest the hand on the knee, shin or ankle.
\textit{Caution:} avoid with a slipped disc, sciatica or after abdominal surgery; do not overdo the side stretch.
\ycbfig{ycb1-fig-trikonasan}

\section{Sitting postures}

\subsection*{Dandasana (staff)}
\textit{Source: NIOS, secondary.} Sit with both legs stretched out and the heels together, palms on the floor beside the buttocks, and the spine, neck and head upright. It is the starting position for several sitting postures: bhadrasana, vakrasana, paschimottanasana and purvottanasana. The relaxed form, shithila dandasana, has the legs a little apart and the trunk leaning slightly back on the hands. The source gives no breath, hold, benefit or caution.
\ycbfig{ycb1-fig-dandasan}

\subsection*{Sukhasana (easy posture)}
\textit{Source: workshop slides.} Sit with the legs crossed comfortably and the hands on the knees or in a mudra; spine upright, shoulders relaxed, eyes closed, attention on the breath. It is credited with calm. With a knee or hip injury sit on a cushion or choose another seat. It is the fallback seat for meditation when padmasana is not possible.
\ycbfig{ycb1-fig-sukhasan}

\subsection*{Padmasana (lotus)}
Sit with the legs straight. Bend one knee and place that foot high on the opposite thigh, then do the same with the other leg. Keep the spine upright and the shoulders relaxed. It is the classic seat for meditation and pranayama. If the full lotus is not possible, use the half form or sukhasana.
\textit{Caution:} avoid with a knee or hip injury, and never force the knees.
\ycbfig{ycb1-fig-padmasan}

\subsection*{Vajrasana (thunderbolt)}
\textit{Source: workshop slides.} Kneel and sit back on the heels, the thighs resting on the calves, hands on the knees palms down, spine straight. It is traditionally the one sitting posture used after a meal (to be checked). It is the starting posture for ardha ushtrasana, shashankasana and mandukasana.
\textit{Caution:} avoid with a knee or ankle injury.
\ycbfig{ycb1-fig-vajrasan}

\subsection*{Bhadrasana (firm or auspicious posture)}
From dandasana, bring the soles of the feet together. Breathing out, clasp the hands over the toes and draw the heels as near the perineum as is comfortable. If the knees stay high, put a soft cushion under them, and stay for some time.
\textit{Caution:} avoid with severe arthritis or sciatica.
\ycbfig{ycb1-fig-bhadrasan}

\subsection*{Mandukasana (frog)}
\textit{Source: CBSE yoga handbook, secondary.} Sit in vajrasana. Make fists with the thumbs tucked inside the fingers and press them against the navel. Breathing out, bend forward with the gaze ahead and hold the breath out in the bent position; then breathe in and come up. Repeat three or four times. It is credited with flexible thighs and legs and with digestion.
\textit{Caution:} avoid with an abdominal injury, back pain or a knee injury.
\ycbfig{ycb1-fig-mandukasan}

\subsection*{Ushtrasana (camel) and ardha ushtrasana (half camel)}
For the half camel, sit on the heels in vajrasana, then rise onto the knees with the hands on the waist, fingers pointing down, elbows and shoulders level. Drop the head back and, breathing in, bend back as far as is comfortable with the thighs upright. Hold for 10 to 30 seconds and return to vajrasana on an in-breath. In the full camel the hands reach the heels, knees about hip width apart, and the hold is with steady breathing; it opens the front of the body and the chest. The half form is the first step.
\textit{Caution:} avoid with hernia, abdominal injury, arthritis, vertigo, back or neck problems and pregnancy.
\ycbfig{ycb1-fig-ushtrasan}

\subsection*{Shashankasana (hare)}
From vajrasana, open the knees wide with the big toes touching and the palms between the knees. Breathing out, stretch the arms forward and bend until the chin rests on the floor, the arms parallel and the eyes looking ahead. Come up on an in-breath and return to vajrasana. It is a natural rest after back bends.
\textit{Caution:} avoid with acute backache; with osteoarthritis of the knee take care or avoid vajrasana.
\ycbfig{ycb1-fig-shashankasan}

\subsection*{Uttana mandukasana (upright frog)}
\textit{Source needed.} The syllabus names it. In one line: a sitting posture in vajrasana with the knees apart, in which the arms are folded behind the head with the elbows as support. No steps are given here until a better source is found; the forward bend shown on old slides is not this posture.
\ycbfig{ycb1-fig-uttanamandukasan}

\subsection*{Paschimottanasana (back stretch)}
\textit{Source: workshop slides.} Sit in dandasana. Breathing in, lengthen the spine; breathing out, bend forward from the hips and take the hands towards the feet. Hold with steady breathing and come up on an in-breath. It stretches the back of the body and the hamstrings and is credited with calm. A stiff person holds the knees or the shins.
\textit{Caution:} avoid with back problems or a hamstring injury.
\ycbfig{ycb1-fig-pashchimottansasan}

\subsection*{Purvottanasana (upward plank, front stretch)}
\textit{Source: Art of Living, secondary.} Sit with the legs out, feet together and back upright, and put the palms on the floor beside the hips, arms straight. Lean back so the hands take the weight, then lift the hips until the body is one straight line, feet on the floor and head tilted back. Breathe in while lifting, breathe steadily in the hold and breathe out while lowering. No hold time is given (other pages say about five breaths), and sources disagree on whether the fingers point forward or back. It strengthens the wrists, arms, shoulders and back, stretches the legs and hips, and is the natural counter to paschimottanasana.
\textit{Caution:} avoid with a wrist or neck injury, or support the head.
\ycbfig{ycb1-fig-purvottansasan}

\subsection*{Vakrasana (spinal twist)}
From dandasana, bend the right leg and place the foot beside the left knee. Breathing out, twist to the right; take the left arm around the right knee and hold the right big toe, or rest the palm beside the right foot. Place the right hand on the floor behind, back straight. Hold for 10 to 30 seconds, release on an out-breath, and repeat on the other side.
\textit{Caution:} avoid with severe back pain, spine or disc disorders, after abdominal surgery and during menstruation.
\ycbfig{ycb1-fig-vakrasan}

\subsection*{Gomukhasana (cow face)}
\textit{Source: Satyananda, a Board-named book; NIOS and CBSE, secondary.} Sit in the folded-leg seat with the right knee directly above the left. Take the left arm behind the back and the right arm over the right shoulder, and try to clasp the fingers behind the back. Bring the raised elbow behind the head so the head rests against the inside of the arm; keep the spine upright and breathe normally. Hold up to two minutes (Satyananda) or one minute (NIOS), then undo and repeat on the other side. It is credited with loose shoulders and neck, an open chest and better posture.
\textit{Caution:} the source gives none; a government yoga page says to avoid it with bleeding piles; take care with a shoulder or knee injury (our advice).
\ycbfig{ycb1-fig-gomukhasan}

\section{Prone postures}

\subsection*{Bhujangasana (cobra) and sarala bhujangasana}
Lie on the stomach, legs together, forehead on the floor. For the easy form, sarala bhujangasana, place the palms and elbows on the floor beside the body and breathe in, lifting the chin and chest to about navel level; then lower. For the full posture, place the palms beside the chest with the elbows raised; breathe in and lift the chin and chest up to the navel region; breathe out and lower. Keep the legs firm so that the lower back is not strained.
\textit{Caution:} after abdominal surgery avoid it for a few months; avoid with hernia or ulcers.
\ycbfig{ycb1-fig-bhujangasan}

\subsection*{Shalabhasana (locust)}
Lie on the stomach with the chin on the floor and the palms up beside the body. Breathe in and raise the straight legs as high as is comfortable. Hold for 10 to 20 seconds with normal breathing, lower on an out-breath, and rest in makarasana. It works better after bhujangasana.
\textit{Caution:} heart patients avoid it; avoid with high blood pressure, peptic ulcer and hernia; with severe low back pain go carefully.
\ycbfig{ycb1-fig-shalabhasan}

\subsection*{Makarasana (crocodile)}
Lie on the stomach with the feet wide apart and pointing out, the arms folded and the forehead on the hands, eyes closed. It is the rest posture between the prone practices.
\textit{Caution:} avoid with low blood pressure, severe heart problems and pregnancy.
\ycbfig{ycb1-fig-makarasan}

\section{Supine postures}

\subsection*{Pavanamuktasana (wind-releasing posture)}
Lie on the back, draw the knees to the chest and clasp the shins below the knees. Breathe out and lift the head until the chin or nose reaches the knees; return the head, and breathe out as the legs go down. Rest in shavasana. Match each movement to the breath.
\textit{Caution:} avoid with an abdominal injury, hernia, sciatica, severe back pain and pregnancy.
\ycbfig{ycb1-fig-pavanamuktasan}

\subsection*{Uttanapadasana and ardha halasana (raised legs)}
\textit{Source: workshop slides.} Lie on the back with the arms beside the body. Breathe in and raise both straight legs to a modest angle, keeping the back and shoulders on the floor. Hold with steady breathing and lower slowly on an out-breath. It strengthens the abdominal muscles. The syllabus lists the two names as separate items; the angle for each is still to be checked.
\textit{Caution:} avoid with back problems.
\ycbfig{ycb1-fig-uttanpadasan}

\subsection*{Setubandhasana (bridge)}
Lie on the back, bend the knees and bring the heels near the buttocks; hold the ankles, knees and feet in one line. Breathe in and lift the buttocks and trunk into a bridge, the shoulders and head staying on the floor; the hands may support the waist. Hold for 10 to 30 seconds, come down on an out-breath, and rest in shavasana.
\textit{Caution:} avoid with ulcers, hernia and late pregnancy.
\ycbfig{ycb1-fig-setubandhasan}

\subsection*{Vipareetakarani (inverted action)}
\textit{Source: Satyananda, a Board-named book; this is his hands-on-hips form, not the legs-up-the-wall form.} Lie flat, legs together, palms down. Raise the straight legs over the body, press on the arms to lift the hips and roll the spine up, then turn the palms up, bend the elbows and cup the hips in the hands near the wrists. Bring the legs to vertical and relax the feet. The trunk ends at about 45 degrees, the weight on the shoulders, neck and elbows, the chin not pressed into the chest. Breathe normally. To come down, lower the legs over the head, put the hands down, lower the spine bone by bone without lifting the head, then lower the straight legs, and rest in shavasana. Beginners hold a few seconds and build up over weeks to 3 to 5 minutes; the knees may bend at first and sore elbows may be padded. It prepares for sarvangasana.
\textit{Caution (as for sarvangasana):} enlarged thyroid, liver or spleen; neck disc problems or spondylitis; high blood pressure or heart problems; weak blood vessels in the eyes; thrombosis; avoid during menstruation and late pregnancy.
\ycbfig{ycb1-fig-vipartikarni}

\subsection*{Saral matsyasana (easy fish)}
\textit{Source: Satyananda, the easy variation of matsyasana with straight legs. The match of this name to that variation is ours and must be confirmed by a certified teacher.} Stretch both legs straight, lean back on the arms, rest the top of the head on the floor and arch the back, with the hands on the thighs or the floor. Return after a while.
\textit{Caution:} heart disease, peptic ulcer, hernia, back conditions, any serious illness and pregnancy; web sources add high or low blood pressure and a delicate neck.
\ycbfig{ycb1-fig-saralmatsyasan}

\subsection*{Shavasana (corpse posture)}
Lie on the back with the arms and legs comfortably apart, palms up, eyes closed. Relax the whole body on purpose, and let the natural breath become slow and rhythmic. Stay until rested. It is the rest posture after the supine practices and the way to end a session of postures.

\section{Counter postures}

\textit{Source: workshop slides, to be checked against the guidebook.} A posture that bends the spine one way is usually followed by rest or a gentle movement the other way. After backward bends such as bhujangasana, shalabhasana or ushtrasana, rest in makarasana or fold forward in shashankasana. After a forward bend such as paschimottanasana, a gentle backward stretch such as purvottanasana can follow.

\ycbfig{ycb1-fig-asanabenefits}
\ycbfig{ycb1-fig-asanabenefits-1}
\ycbfig{ycb1-fig-asanabenefits-2}
\ycbfig{ycb1-fig-asanabenefits-3}

\begin{recall}
\begin{itemize}
\item For every asana: name, start, breath, hold, way out, one benefit, one caution.
\item Typical protocol hold: 10 to 30 seconds with normal breathing.
\item Rest postures: makarasana between prone practices, shavasana after supine ones.
\item Starting seats: dandasana for bhadrasana, vakrasana, paschimottanasana, purvottanasana; vajrasana for ardha ushtrasana, shashankasana, mandukasana.
\item Benefits are traditional credits, not cures; a caution means modify or skip.
\end{itemize}
\end{recall}

\subsection*{Worked examples}
\begin{itemize}
\item \textbf{A student with a stiff back cannot bend forward in padahastasana.} Let the hands rest on the thighs or a chair seat, bend only as far as the back allows, and explain that the aim is the stretch of the back of the legs, not touching the floor.
\item \textbf{A pregnant student joins a mixed class.} Leave out padahastasana, ushtrasana, makarasana, pavanamuktasana and the inverted posture; offer seated and gentle standing postures, and ask the student to get a doctor's advice.
\item \textbf{Model answer: teach trikonasana in one minute.} Feet apart, front foot out; arms at shoulder height; breathe out and bend sideways, hand behind the foot or on the shin; top arm up, look at the fingertips; hold 10 to 30 seconds; come up on an in-breath; other side. Credited with a side stretch of the trunk; avoid with a slipped disc, sciatica or after abdominal surgery.
\end{itemize}

\subsection*{Check yourself}
\begin{enumerate}
\item How long is the hold in trikonasana in the protocol?
\item Which posture is the rest between prone practices?
\item Name two conditions in which padahastasana is avoided.
\item Which posture is the starting position for bhadrasana, vakrasana and paschimottanasana?
\item Which sitting postures start from vajrasana?
\item Which seat is the classic one for meditation and pranayama, with each foot on the opposite thigh?
\item Describe the arm position of gomukhasana.
\item What follows a series of backward bends as a rest?
\item In the standing twist of kati chakrasana, how many rounds are done?
\item What is the end position of the trunk in vipareetakarani, in the form given here?
\end{enumerate}

\subsection*{Answers}
\begin{enumerate}
\item About 10 to 30 seconds with normal breathing.
\item Makarasana.
\item Any two of: heart or back problems, hernia, ulcers, vertigo, high myopia, pregnancy.
\item Dandasana.
\item Ardha ushtrasana, shashankasana and mandukasana.
\item Padmasana.
\item One arm over the shoulder and the other behind the back, the fingers clasped behind.
\item Makarasana or shashankasana.
\item Five to ten.
\item At about 45 degrees, the hands cupping the hips.
\end{enumerate}

\paragraph{Sources.} Common Yoga Protocol, Ministry of AYUSH (the protocol postures and hold times); the Board syllabus for Level 1 (the list and grouping); Swami Satyananda Saraswati, \emph{Asana Pranayama Mudra Bandha} (kati chakrasana, gomukhasana, vipareetakarani, the easy fish); National Institute of Open Schooling, Yoga (dandasana, gomukhasana); CBSE Class XI yoga handbook (mandukasana); Art of Living (purvottanasana); the official guidebook for Level I (2016, older scheme) as a coverage check. Benefits credited by these sources are given as traditional credits; their cure claims are not repeated.
